%% subject_reduction_v8.tex
%% Subject Reduction (Preservation) + Progress for F_TG
%% Appendix section -- input from appendix_v6.tex

%% Semantic brackets (stmaryrd not available; define locally)
\providecommand{\llbracket}{[\![}
\providecommand{\rrbracket}{]\!]}

\section{Subject Reduction and Progress for
         \texorpdfstring{$\mathcal{F}_{\mathrm{TG}}$}{F-TG}}
\label{app:subject-reduction}

This section gives a self-contained small-step operational semantics for the
statement fragment $\mathcal{F}_{\mathrm{TG}}$, establishes Preservation
(\Cref{thm:preservation}) and Progress (\Cref{thm:progress}) as formal theorems,
derives the classical Soundness corollary, and provides full proofs.
Every operator typing rule cited is from \Cref{sec:calculus,app:typing-rules};
no rule is invented here.

%% ─────────────────────────────────────────────────────────────────────────────
\subsection{Runtime Configurations}
\label{app:sos-configs}

A \emph{runtime value} $v$ is a record
$(\mathit{sh},\tau,\mathit{rg},\mathit{ot})$ where
$\mathit{sh}\in\mathbb{Z}_{>0}^{*}$ is the concrete shape tuple,
$\tau\in\mathsf{DType}$, $\mathit{rg}\in\mathbb{B}$ is
\texttt{requires\_grad}, and $\mathit{ot}\in\mathbb{B}$ is
\texttt{is\_on\_tape}.
A \emph{heap} $\sigma:\mathsf{Var}\rightharpoonup\mathsf{Val}$ maps variable
names to runtime values; $\sigma[x\mapsto v]$ denotes functional update.
The \emph{grad context} $G\in\{\mathsf{on},\mathsf{off}\}$ records whether
the cursor is inside \texttt{torch.no\_grad()}, matching the flag in
\Cref{sec:calculus}.

\begin{definition}[Heap satisfaction, $\sigma\models\Gamma$]
\label{def:heap-satisfaction}
$\sigma\models\Gamma$ iff for every binding
$x:\mathsf{Tensor}[\tau,\varphi_s,\varphi_g]\in\Gamma$,
writing $\sigma(x)=(\mathit{sh},\tau',\mathit{rg},\mathit{ot})$:
\begin{enumerate}
  \item $\tau'=\tau$;
  \item $\varphi_s[\mathit{sh}/\bar d]$ holds (concrete instantiation
        of the shape predicate over $\mathcal{T}_{\mathsf{shape}}$);
  \item if $\varphi_g=\mathsf{has\_grad}$ then $\mathit{rg}\wedge\mathit{ot}=\mathsf{true}$;
        if $\varphi_g=\mathsf{no\_grad}$ then $\mathit{rg}\wedge\mathit{ot}=\mathsf{false}$;
        if $\varphi_g=\top$ the condition is vacuous.
\end{enumerate}
For $\mathsf{unit}$ bindings and static-integer bindings the condition holds
trivially.
\end{definition}

We extend the statement grammar of \Cref{sec:fragment} with two meta-forms used
only in reduction rules:
\[
  s \;::=\; \cdots \;\mid\; s_1\,;\,s_2 \;\mid\; \mathbf{skip}
\]
A \emph{program configuration} is $\langle s,\sigma,G\rangle$;
$\langle\mathbf{return}\;v,\sigma,G\rangle$ is \emph{terminal}
(a value form); $\langle\textsc{Error},\sigma,G\rangle$ is the
\emph{error sink}.

%% ─────────────────────────────────────────────────────────────────────────────
\subsection{Small-Step Reduction Rules}
\label{app:sos-rules}

\paragraph{Expression reduction
  $\langle e,\sigma,G\rangle\to_e\langle e',\sigma,G\rangle$.}

{\small
\begin{mathpar}
\inferrule*[lab=\textsc{E-Var}]
  {x\in\mathrm{dom}(\sigma)}
  {\langle x,\sigma,G\rangle\to_e\langle\sigma(x),\sigma,G\rangle}
\and
\inferrule*[lab=\textsc{E-Op}]
  {\sigma(\bar x)=\bar v \quad \mathsf{prec}(\mathit{op},\bar v)\;\text{holds}}
  {\langle x_1.\mathit{op}(\bar x_{2..}),\sigma,G\rangle
    \to_e \langle\llbracket\mathit{op}\rrbracket(\bar v,G),\sigma,G\rangle}
\and
\inferrule*[lab=\textsc{E-OpErr}]
  {\sigma(\bar x)=\bar v \quad \mathsf{prec}(\mathit{op},\bar v)\;\text{fails}}
  {\langle x_1.\mathit{op}(\bar x_{2..}),\sigma,G\rangle
    \to_e \langle\textsc{Error},\sigma,G\rangle}
\and
\inferrule*[lab=\textsc{E-NgBody}]
  {\langle e,\sigma,\mathsf{off}\rangle\to_e\langle e',\sigma,\mathsf{off}\rangle}
  {\langle\mathtt{no\_grad}(e),\sigma,G\rangle
    \to_e \langle\mathtt{no\_grad}(e'),\sigma,G\rangle}
\and
\inferrule*[lab=\textsc{E-NgVal}]
  {v=(\mathit{sh},\tau,\mathit{rg},\mathit{ot}) \quad
   v'=(\mathit{sh},\tau,\mathit{rg},\mathsf{false})}
  {\langle\mathtt{no\_grad}(v),\sigma,G\rangle\to_e\langle v',\sigma,G\rangle}
\end{mathpar}}

\noindent
Here $\mathsf{prec}(\mathit{op},\bar v)$ is the shape-compatibility predicate
for $\mathit{op}$ (e.g.\ equal contracting dimensions for \texttt{matmul};
divisibility for \texttt{view($-1$)}), and $\llbracket\mathit{op}\rrbracket(\bar v,G)$
is the concrete result whose shape is computed by the same transfer function as
the typing rule.
\textsc{E-NgBody} / \textsc{E-NgVal} evaluate a \texttt{no\_grad}-wrapped
expression with $G=\mathsf{off}$ and strip the tape bit on exit, matching
\textsc{T-NoGradCtx}.
Tuple construction, module calls $C(\bar x)$, and attribute reads
$\mathit{self}.\mathit{attr}$ reduce component-wise via standard eval-context
rules (omitted).

\paragraph{Statement reduction
  $\langle s,\sigma,G\rangle\to\langle s',\sigma',G'\rangle$.}

{\small
\begin{mathpar}
\inferrule*[lab=\textsc{E-EvalCtx}]
  {\langle e,\sigma,G\rangle\to_e\langle e',\sigma,G\rangle}
  {\langle x=e\,;\,s,\sigma,G\rangle\to\langle x=e'\,;\,s,\sigma,G\rangle}
\and
\inferrule*[lab=\textsc{E-Bind}]
  {v\;\text{is a value}}
  {\langle x=v\,;\,s,\sigma,G\rangle\to\langle s,\sigma[x\mapsto v],G\rangle}
\and
\inferrule*[lab=\textsc{E-EvalErr}]
  {\langle e,\sigma,G\rangle\to_e\langle\textsc{Error},\sigma,G\rangle}
  {\langle x=e\,;\,s,\sigma,G\rangle\to\langle\textsc{Error},\sigma,G\rangle}
\and
\inferrule*[lab=\textsc{E-IfT}]
  {c\;\text{evaluates to}\;\mathsf{true}\;\text{in}\;\sigma}
  {\langle\mathbf{if}^\dagger\,c:\,s_1\,\mathbf{else}\,s_2\,;\,s,\sigma,G\rangle
   \to\langle s_1\,;\,s,\sigma,G\rangle}
\and
\inferrule*[lab=\textsc{E-IfF}]
  {c\;\text{evaluates to}\;\mathsf{false}\;\text{in}\;\sigma}
  {\langle\mathbf{if}^\dagger\,c:\,s_1\,\mathbf{else}\,s_2\,;\,s,\sigma,G\rangle
   \to\langle s_2\,;\,s,\sigma,G\rangle}
\and
\inferrule*[lab=\textsc{E-ForZ}]
  { }
  {\langle\mathbf{for}^\dagger\,x\,\mathbf{in}\,\mathbf{range}(0):\,s_b\,;\,s,\sigma,G\rangle
   \to\langle s,\sigma,G\rangle}
\and
\inferrule*[lab=\textsc{E-ForS}]
  {n>0}
  {\langle\mathbf{for}^\dagger\,x\,\mathbf{in}\,\mathbf{range}(n):\,s_b\,;\,s,\sigma,G\rangle
   \to\langle s_b[0/x]\,;\,
         \mathbf{for}^\dagger\,x\,\mathbf{in}\,\mathbf{range}(n{-}1):\,s_b[(x{+}1)/x]
         \,;\,s,\sigma,G\rangle}
\and
\inferrule*[lab=\textsc{E-Skip}]
  { }
  {\langle\mathbf{skip}\,;\,s,\sigma,G\rangle\to\langle s,\sigma,G\rangle}
\end{mathpar}}

\noindent
In \textsc{E-ForS}, $s_b[(x{+}1)/x]$ denotes the substitution of $x{+}1$ for
every free occurrence of $x$ in $s_b$; iteration $i$ of the residual body
receives index $i{+}1$, unrolling correctly to iterations $1,\dots,n{-}1$.
Branch determinism (\textsc{E-IfT}/\textsc{E-IfF}) holds by the
shape-determinism restriction: $c$ is a static integer expression, and
$\sigma\models\Gamma$ instantiates all shape variables, so $c$ has a unique
Boolean value.
Loop termination follows because the bound $n$ strictly decreases in
\textsc{E-ForS} and is non-negative (it is a natural number derivable from
shape refinements).

%% ─────────────────────────────────────────────────────────────────────────────
\subsection{Statement Typing}
\label{app:stmt-typing}

The expression typing judgment $\Gamma\vdash e:T$ is from \Cref{sec:calculus}.
We extend it to statements with $\Gamma\vdash s:\tau$, meaning that $s$
terminates with a value of type $\tau$ in context $\Gamma$.

{\small
\begin{mathpar}
\inferrule*[lab=\textsc{T-Assign}]
  {\Gamma\vdash e:T \quad \Gamma,x{:}T\vdash s:\tau}
  {\Gamma\vdash (x=e\,;\,s):\tau}
\and
\inferrule*[lab=\textsc{T-Return}]
  {\Gamma\vdash e:\tau}
  {\Gamma\vdash\mathbf{return}\;e:\tau}
\and
\inferrule*[lab=\textsc{T-If}]
  {\Gamma\vdash s_1:\tau \quad \Gamma\vdash s_2:\tau \quad
   c\;\text{static in}\;\Gamma}
  {\Gamma\vdash (\mathbf{if}^\dagger\,c:\,s_1\,\mathbf{else}\,s_2):\tau}
\and
\inferrule*[lab=\textsc{T-For}]
  {\Gamma,x{:}\mathsf{int}[0..n)\vdash s_b:\mathsf{unit} \quad
   n\;\text{static in}\;\Gamma}
  {\Gamma\vdash (\mathbf{for}^\dagger\,x\,\mathbf{in}\,\mathbf{range}(n):\,s_b):
   \mathsf{unit}}
\and
\inferrule*[lab=\textsc{T-Seq}]
  {\Gamma\vdash s_1:\mathsf{unit} \quad \Gamma'\vdash s_2:\tau \quad
   \Gamma'\supseteq\Gamma}
  {\Gamma\vdash (s_1\,;\,s_2):\tau}
\and
\inferrule*[lab=\textsc{T-Skip}]
  { }
  {\Gamma\vdash\mathbf{skip}:\mathsf{unit}}
\end{mathpar}}

\noindent
\textsc{T-If} requires both branches to return the \emph{same} type $\tau$,
using shape-determinism to avoid a lattice join on the heap.
\textsc{T-Assign} uniquely determines $T$ via one of the expression rules
(\textsc{T-Matmul}, \textsc{T-View($-1$)}, etc.) from \Cref{sec:calculus,app:typing-rules}.

%% ─────────────────────────────────────────────────────────────────────────────
\subsection{Preliminary Lemmas}

\begin{lemma}[Canonical Forms]\label{lem:canonical}
If $\Gamma\vdash x:\mathsf{Tensor}[\tau,\varphi_s,\varphi_g]$,
$x\in\mathrm{dom}(\sigma)$, and $\sigma\models\Gamma$, then
writing $\sigma(x)=(\mathit{sh},\tau',\mathit{rg},\mathit{ot})$:
$\tau'=\tau$, $\varphi_s[\mathit{sh}/\bar d]$ holds, and the grad condition
of \Cref{def:heap-satisfaction}(iii) is satisfied.
\end{lemma}
\begin{proof}
Immediate from \Cref{def:heap-satisfaction}.
\end{proof}

\begin{lemma}[Substitution]\label{lem:substitution}
If $\Gamma,x{:}S\vdash s:\tau$ and $\Gamma\vdash v:S$,
then $\Gamma\vdash s[v/x]:\tau$.
\end{lemma}
\begin{proof}[Sketch]
Standard structural induction on $s$.
Each use of $x$ in $s$ is replaced by $v$; the unique type $S$ is preserved at
every use site.
Shape predicates mentioning $x$ are instantiated pointwise in
$\mathcal{T}_{\mathsf{shape}}$.
\end{proof}

\begin{lemma}[Reduction Lemma]\label{lem:reduction}
If $\Gamma\vdash e:T$, $\sigma\models\Gamma$, and
$\langle e,\sigma,G\rangle\to_e\langle v,\sigma,G\rangle$ via
\textsc{E-Op} (non-error), then $v$ satisfies type $T$ in $\sigma$.
\end{lemma}
\begin{proof}
By case analysis on the expression typing rule for $e$.

\emph{Case \textsc{E-Var} / \textsc{T-Var}.}
$e=x$, $v=\sigma(x)$.
By \Cref{lem:canonical}, $v$ satisfies $T$.

\emph{Case \textsc{T-Matmul}.}
Premises: $a:\mathsf{Tensor}[\tau,(\bar B,m,k),\varphi_a]$,
$b:\mathsf{Tensor}[\tau,(\bar B',k,n),\varphi_b]$,
$\Gamma\models\mathrm{broadcast}(\bar B,\bar B')=\bar C$.
By \Cref{lem:canonical}, the concrete shapes satisfy these conditions.
\textsc{E-Op} fires because $\mathsf{prec}(\mathtt{matmul},\bar v)$
(equal contracting dimension) holds.
PyTorch returns $(\bar C_\sigma,m_\sigma,n_\sigma)$ where $\bar C_\sigma$
is the concrete broadcast of $\bar B_\sigma$ and $\bar B'_\sigma$.
The type $\mathsf{Tensor}[\tau,(\bar C,m,n),g(a,b)]$ is satisfied:
$\bar C_\sigma$ concretely instantiates $\bar C$, and
$g(a,b)=\bigsqcup(\varphi_a^\mathsf{g},\varphi_b^\mathsf{g})$ matches
PyTorch's autograd contract
(output $\mathit{ot}=\mathsf{true}$ iff some input has $\mathit{ot}=\mathsf{true}$
and $G=\mathsf{on}$).

\emph{Case \textsc{T-View($-1$)}.}
Premises: $\Gamma\models Q\mid P$ where $P=\prod_i d_i$, $Q=\prod_{i\neq j}s_i$.
\textsc{E-Op} fires only when $Q_\sigma\mid P_\sigma$; the inferred
dimension $P_\sigma/Q_\sigma$ is the concrete instantiation of $P/Q$.
Shape and grad are preserved (\texttt{view} is metadata-only).

\emph{Case \textsc{T-Detach}.}
The result has $\mathit{ot}=\mathsf{false}$ by PyTorch's \texttt{detach}
contract, satisfying $\mathsf{no\_grad}$.

\emph{Case \textsc{T-NoGradCtx}.}
By \textsc{E-NgVal}, the output value has $\mathit{ot}=\mathsf{false}$
regardless of its origin, satisfying $\mathsf{no\_grad}$ as asserted.

\emph{Case \textsc{T-Conv2d}.}
The output dimensions $H',W'$ are deterministic arithmetic functions of
$H,W$ and the convolution hyperparameters (stride, padding, dilation, kernel);
PyTorch computes exactly the same formula.
Grad: as in \textsc{T-Matmul}.

\emph{Case \textsc{T-Cat}.}
The concatenated dimension is $\sum_i d_k^{(i)}$; the concrete value is
$\sum_i\sigma(d_k^{(i)})$, instantiating the symbolic type.

\emph{Case \textsc{T-Split}.}
\textsc{E-Op} fires when $s_\sigma\mid d_{k,\sigma}$; each output chunk has
concrete size $s_\sigma=\sigma(s)$, matching $s$.

\emph{Case \textsc{T-SDPA}.}
Premise $T=S$ is checked by $\mathsf{prec}$; output shape $(B,H,T,D)$ matches.

\emph{Case \textsc{T-LayerNorm}.}
Output shape equals input shape (only rank-matching is required);
grad as in \textsc{T-Matmul}.

\emph{Case \textsc{T-Backward} / \textsc{T-OptStep}.}
Result type is $\mathsf{unit}$; the condition is vacuous.

\emph{Lean-audited operators}
(\Cref{tab:handler-soundness}, 28 operators including \textsc{T-Linear},
\textsc{T-Bmm}, \textsc{T-Permute}, \textsc{T-Expand}, \textsc{T-Repeat},
\textsc{T-Stack}, \textsc{T-Chunk}, \textsc{T-Unbind},
\textsc{T-Gather}, \textsc{T-Scatter}, \textsc{T-IndexSelect},
\textsc{T-Narrow}, \textsc{T-Embed}, \textsc{T-RmsNorm}).
Each Lean definition formalises: for any input satisfying the rule's premises,
the rule's output shape formula equals PyTorch's runtime output---this is exactly
the statement of the present lemma for each such operator, and the Lean proof
discharges it.

\emph{Pen-and-paper operators} (\textsc{T-Broadcast}, \textsc{T-Reduce},
\textsc{T-Einsum}).
Stated and proved in the pen-and-paper note of
\Cref{tab:handler-soundness}; shape transfer is exact arithmetic and
is verified by the same argument as \textsc{T-Cat} above.

\emph{Tested-only operators}: outside the scope of this theorem
(cf.\ the boundary note in Theorem~\ref{thm:fragment-soundness}).
\end{proof}

%% ─────────────────────────────────────────────────────────────────────────────
\subsection{Preservation (Subject Reduction)}

\begin{theorem}[Preservation]\label{thm:preservation}
If $\Gamma\vdash s:\tau$, $\sigma\models\Gamma$, and
$\langle s,\sigma,G\rangle\to\langle s',\sigma',G'\rangle$
is a non-error step,
then there exists $\Gamma'\supseteq\Gamma$ such that
$\sigma'\models\Gamma'$ and $\Gamma'\vdash s':\tau$.
\end{theorem}
\begin{proof}
By case analysis on the reduction rule applied.

\paragraph{Case \textsc{E-EvalCtx}.}
$s = x=e\,;\,s_0$, step has $e\to_e e'$ (non-error),
$s'=x=e'\,;\,s_0$, $\sigma'=\sigma$, $G'=G$.
By \textsc{T-Assign}, $\Gamma\vdash e:T$ for some $T$.
By \Cref{lem:reduction} (applied to the sub-expression step),
the partially-reduced $e'$ still types at $T$ under $\Gamma$
(expression reduction does not change $\sigma$, so $\sigma\models\Gamma$ is
preserved).
Setting $\Gamma'=\Gamma$: $\sigma'\models\Gamma'$ and
$\Gamma\vdash(x=e'\,;\,s_0):\tau$ by \textsc{T-Assign}.

\paragraph{Case \textsc{E-Bind}.}
$s=x=v\,;\,s_0$, $s'=s_0$, $\sigma'=\sigma[x\mapsto v]$, $G'=G$.
By \textsc{T-Assign}: $\Gamma\vdash v:T$ (already a value at type $T$
because the preceding \textsc{E-EvalCtx} steps preserved typing, by the
previous case) and $\Gamma,x{:}T\vdash s_0:\tau$.
Since $v$ satisfies $T$ (by \Cref{lem:reduction} on the completed
expression-evaluation sub-sequence), $\sigma[x\mapsto v]\models\Gamma,x{:}T$:
all old bindings are unchanged and $x{:}T$ is freshly satisfied.
Set $\Gamma'=\Gamma,x{:}T$; then $\sigma'\models\Gamma'$ and
$\Gamma'\vdash s_0:\tau$.

\paragraph{Case \textsc{E-IfT}.}
$s=\mathbf{if}^\dagger\,c:\,s_1\,\mathbf{else}\,s_2\,;\,s_0$.
By \textsc{T-If}, $\Gamma\vdash s_1:\tau$.
Step gives $s'=s_1\,;\,s_0$; by \textsc{T-Seq},
$\Gamma\vdash s_1\,;\,s_0:\tau$.
$\sigma'=\sigma$, $G'=G$, $\Gamma'=\Gamma$.

\paragraph{Case \textsc{E-IfF}.}
Symmetric; $\Gamma\vdash s_2:\tau$ from \textsc{T-If}.

\paragraph{Case \textsc{E-ForZ}.}
$s=\mathbf{for}^\dagger\,x\,\mathbf{in}\,\mathbf{range}(0):\,s_b\,;\,s_0$.
Step gives $s'=s_0$.
By \textsc{T-For}, the for-loop has type $\mathsf{unit}$; by \textsc{T-Seq},
$\Gamma\vdash s_0:\tau$.  $\Gamma'=\Gamma$.

\paragraph{Case \textsc{E-ForS}.}
$s=\mathbf{for}^\dagger\,x\,\mathbf{in}\,\mathbf{range}(n):\,s_b\,;\,s_0$,
$n>0$.
By \textsc{T-For}: $\Gamma,x{:}\mathsf{int}[0..n)\vdash s_b:\mathsf{unit}$;
the step produces
$s'=s_b[0/x]\,;\,
   \mathbf{for}^\dagger\,x\,\mathbf{in}\,\mathbf{range}(n{-}1):\,
   s_b[(x{+}1)/x]\,;\,s_0$.
\begin{enumerate}
\item $\Gamma\vdash s_b[0/x]:\mathsf{unit}$:
  since $n>0$, $0\in[0..n)$, so \Cref{lem:substitution} applies.
\item $\Gamma,x{:}\mathsf{int}[0..n{-}1)\vdash s_b[(x{+}1)/x]:\mathsf{unit}$:
  in the residual body $x$ ranges over $[0..n{-}1)$ and $x{+}1$ ranges over
  $[1..n)\subseteq[0..n)$; \Cref{lem:substitution} (with term
  $x{+}1:\mathsf{int}[1..n)$) gives this.
\item By \textsc{T-For} with bound $n{-}1$:
  $\Gamma\vdash
   \mathbf{for}^\dagger\,x\,\mathbf{in}\,\mathbf{range}(n{-}1):\,
   s_b[(x{+}1)/x]:\mathsf{unit}$.
\item By \textsc{T-Seq} applied twice: $\Gamma\vdash s':\tau$.
\end{enumerate}
$\Gamma'=\Gamma$, $\sigma'=\sigma$.

\paragraph{Case \textsc{E-Skip}.}
$s=\mathbf{skip}\,;\,s_0\to s_0$.
By \textsc{T-Seq}, $\Gamma\vdash s_0:\tau$.
$\Gamma'=\Gamma$, $\sigma'=\sigma$.

\smallskip
All cases produce $\Gamma'\supseteq\Gamma$ with
$\sigma'\models\Gamma'$ and $\Gamma'\vdash s':\tau$.
\end{proof}

%% ─────────────────────────────────────────────────────────────────────────────
\subsection{Progress}

\begin{theorem}[Progress]\label{thm:progress}
If $\Gamma\vdash s:\tau$ and $\sigma\models\Gamma$, then one of:
\begin{enumerate}
  \item $s=\mathbf{return}\;v$ is terminal (value form);
  \item there exists $\langle s',\sigma',G'\rangle$ such that
        $\langle s,\sigma,G\rangle\to\langle s',\sigma',G'\rangle$; or
  \item $\langle s,\sigma,G\rangle\to\langle\textsc{Error},\sigma,G\rangle$,
        and this arises only when $\mathsf{prec}(\mathit{op},\bar v)$
        fails for some operator $\mathit{op}$ in $s$---equivalently,
        when the corresponding typing-rule premise
        (\textsc{T-Matmul}'s $k$-equality,
        \textsc{T-View($-1$)}'s divisibility, etc.)
        is not satisfied by the concrete shapes in $\sigma$.
\end{enumerate}
\end{theorem}
\begin{proof}
By structural induction on the typing derivation $\Gamma\vdash s:\tau$.

\emph{Case \textsc{T-Return}.}
$s=\mathbf{return}\;e$.
If $e$ is a value, clause~(a).
Otherwise, locate the innermost non-value sub-expression of $e$;
by the grammar of $\mathcal{F}_{\mathrm{TG}}$, it is of the form
$x_1.\mathit{op}(\bar x_{2..})$ with all $x_i\in\mathrm{dom}(\sigma)$.
If $\mathsf{prec}(\mathit{op},\sigma(\bar x))$ holds, \textsc{E-Op} fires and
the result propagates via \textsc{E-EvalCtx} to produce a step (clause b);
if it fails, \textsc{E-OpErr} fires followed by \textsc{E-EvalErr}
(clause c with the named precondition failure).

\emph{Case \textsc{T-Assign}.}
$s=x=e\,;\,s_0$.
Identical argument on $e$; if $e$ is already a value,
\textsc{E-Bind} fires (clause b).

\emph{Case \textsc{T-If}.}
The shape-determinism restriction ensures $c$ is a static integer expression.
Since $\sigma\models\Gamma$, all shape variables in $c$ are instantiated
and $c$ evaluates to a unique Boolean: \textsc{E-IfT} or \textsc{E-IfF}
applies.  Clause~(b).

\emph{Case \textsc{T-For}.}
The loop bound $n$ is static in $\Gamma$ and hence determined by $\sigma$;
$n\geq 0$ since it is a natural number derivable from shape refinements.
If $n=0$, \textsc{E-ForZ} applies; if $n>0$, \textsc{E-ForS} applies.
Clause~(b).

\emph{Case \textsc{T-Seq}.}
$s=s_1\,;\,s_2$.
If $s_1=\mathbf{skip}$, \textsc{E-Skip} applies (clause b).
Otherwise the IH on $s_1$ gives clause~(b) (a step for $s_1$,
propagated to $s_1\,;\,s_2$ by the eval-context rule for sequences)
or clause~(c) (an error propagated by \textsc{E-EvalErr}).

\smallskip\noindent
No other stuck states arise within $\mathcal{F}_{\mathrm{TG}}$:
data-dependent branches are excluded by definition;
dynamic dispatch outside $\mathrm{Cat}$ triggers \textsc{Abstain} at the
analysis phase and does not reach these semantics.
\end{proof}

%% ─────────────────────────────────────────────────────────────────────────────
\subsection{Soundness Corollary}

\begin{corollary}[Soundness, classical form]\label{cor:soundness}
Let $s\in\mathcal{F}_{\mathrm{TG}}$ with $\Gamma\vdash s:\tau$ and
$\sigma\models\Gamma$.
If every operator $\mathit{op}$ reachable during reduction of
$\langle s,\sigma,G\rangle$ satisfies $\mathsf{prec}(\mathit{op},\bar v)$
in $\sigma$, then no reduction sequence raises
\textsc{Shape-Mismatch-Error}.
\end{corollary}
\begin{proof}
Under the hypothesis, \Cref{thm:progress}(c) never applies.
\Cref{thm:preservation} guarantees that each non-error step yields a
well-typed configuration $(\Gamma'\supseteq\Gamma,\,s',\sigma')$ with
$\sigma'\models\Gamma'$, so the precondition hypothesis is maintained at every
subsequent step.
By induction on the length of any finite reduction prefix, the error sink is
never reached.
\end{proof}

\paragraph{Relationship to Theorem~\ref{thm:soundness}.}
Theorem~\ref{thm:soundness} (soundness of proof-grade verdicts) is the
\emph{abstract-interpretation specialisation} of
\Cref{thm:preservation,thm:progress} to the verdict lattice.
Concretely: \textsc{Verified} means Z3 has symbolically discharged
$\mathsf{prec}(\mathit{op},\cdot)$ for all operators under $\Gamma$, so the
precondition hypothesis of \Cref{cor:soundness} holds for every
$\sigma\models\Gamma$, yielding Theorem~\ref{thm:soundness}(i).
\textsc{Refuted-Proof} witnesses a concrete $\sigma\models\Gamma$ in which some
$\mathsf{prec}$ fails, so \Cref{thm:progress}(c) fires at the witnessed step,
giving Theorem~\ref{thm:soundness}(ii).
Hence Theorem~\ref{thm:soundness} is exactly the pair
$(\Cref{thm:preservation},\Cref{thm:progress})$ run through the Z3 abstraction
layer; the verdict lattice is the coarsened output domain.
